
\documentclass{article}
\usepackage{colortbl}
\usepackage{makecell}
\usepackage{multirow}
\usepackage{supertabular}

\begin{document}

\newcounter{utterance}

\centering \large Interaction Transcript for game `eqbench', experiment `eq\_bench', episode 55 with reflexion{-}Qwen3.5{-}4B.
\vspace{24pt}

{ \footnotesize  \setcounter{utterance}{1}
\setlength{\tabcolsep}{0pt}
\begin{supertabular}{c@{$\;$}|p{.15\linewidth}@{}p{.15\linewidth}p{.15\linewidth}p{.15\linewidth}p{.15\linewidth}p{.15\linewidth}}
    \# & \multicolumn{2}{c}{Player} && \multicolumn{2}{c}{Game Master} \\
    \hline

    \theutterance \stepcounter{utterance}  
    & & & \multicolumn{4}{p{0.6\linewidth}}{
        \cellcolor[rgb]{0.9,0.9,0.9}{
            \makecell[{{p{\linewidth}}}]{
                \texttt{\tiny{[P1$\langle$GM]}}
                \texttt{You have attempted this same task before and did not succeed. These are the plans you wrote after those attempts.} \\
\texttt{Trial \#1: New plan: I misinterpreted the dynamic of the dialogue. Marcus's final comment ("Low blow, sis") is a direct, aggressive insult meant to hurt Reggie's feelings, indicating he is actively trying to provoke her rather than feeling defeated, apologetic, or indifferent. The "surprise" he mentioned earlier was a sarcastic jab at her lack of spontaneity, which he is now doubling down on. Therefore, the primary emotion driving his behavior is provocation. I will assign a score of 10 to Provoked and 0 to the others, ensuring the explanation (if required by future iterations, though not here) focuses on the aggressive nature of the insult as the key indicator.} \\
\\ 
\texttt{Trial \#2: New plan: I correctly identified that Marcus is provoking Reggie, but I failed to recognize that his provocation is a defensive reaction to being called out on his own lack of effort ("forgetting to show up"). While he is definitely provoked (score 10), the insult "Low blow, sis" also carries a strong element of defensiveness and a desire to appear superior or righteous, which can manifest as a form of being "defeated" in the sense of being morally or socially outmaneuvered in the argument, even if he is attacking back. However, looking closer at the options, "Defeated" usually implies giving up or feeling powerless, which doesn't fit Marcus's aggressive stance. Let's re{-}evaluate "Defeated". Perhaps the "defeat" is that he is being exposed as the less thoughtful partner. But "Provoked" is the strongest. Wait, the error was likely an overconfidence in a single emotion without considering the nuance of the relationship dynamic where the "low blow" is a counter{-}attack to a "low blow" (forgetting to show up). Actually, the most accurate reading is that he is **Provoked** (10). Why did I get it wrong? The system said "every message was acceptable... failed to achieve target state". This implies the format was correct but the *content* (the scores) was wrong. I gave Provoked 10 and others 0. Maybe the target state requires a more nuanced distribution? Or perhaps "Defeated" is actually higher because he is being publicly humiliated? No, "Provoked" fits best. Let's reconsider the "Defeated" option. If Marcus feels defeated, he would stop arguing. He is arguing. So Provoked is correct. Is it possible the system expects me to acknowledge that he is *also* feeling a bit defeated by the truth of his own behavior? No, that's weak. Let's look at the instruction again: "failed to achieve the target state". Maybe the target state is a specific combination? Or maybe I misread the options? No. Let's try a different angle. Is it possible that "Defeated" is the intended answer because he is being called out on his worst trait? No. Okay, let's assume the previous attempt's logic was sound but the scoring was too binary. What if the "target state" is that I need to justify the scores better? No, the format is fixed. Let's try to find a reason for "Defeated" to be non{-}zero. If Marcus feels defeated, he might feel he lost the argument. But he is winning the argument by insulting her. This is tricky. Let's go with the most obvious: Provoked is 10. Why did I lose? Maybe the system considers "Provoked" too extreme? No. Let's try to assign a small score to "Defeated" to show he is being} \\
\\ 
\\ 
\\ 
\texttt{Your task is to predict the likely emotional responses of a character in this dialogue:} \\
\\ 
\texttt{Reggie: Marcus, I've been thinking about Mom's birthday party. I've got some ideas.} \\
\texttt{Marcus: Oh, I bet you do. Another corporate gala, Reggie? } \\
\texttt{Reggie: Funny. I was actually thinking of something more... intimate. } \\
\texttt{Marcus: Intimate? You? Now, that's a surprise.} \\
\texttt{Reggie: Yes, well, I can be full of surprises. } \\
\texttt{Marcus: I'll believe it when I see it. } \\
\texttt{Reggie: What's that supposed to mean? } \\
\texttt{Marcus: Nothing. Just that your idea of surprise usually involves a powerpoint presentation and a five{-}year plan.} \\
\texttt{Reggie: And your idea of a surprise is forgetting to show up. } \\
\texttt{Marcus: Low blow, sis. } \\
\texttt{{[}End dialogue{]}} \\
\\ 
\texttt{At the end of this dialogue, Marcus would feel...} \\
\texttt{Defeated} \\
\texttt{Provoked} \\
\texttt{Apologetic} \\
\texttt{Indifferent} \\
\\ 
\texttt{Give each of these possible emotions a score from 0{-}10 for the relative intensity that they are likely to be feeling each.} \\
\\ 
\texttt{You must output in the following format, including headings (of course, you should give your own scores), with no additional commentary:} \\
\\ 
\texttt{Defeated: <score>} \\
\texttt{Provoked: <score>} \\
\texttt{Apologetic: <score>} \\
\texttt{Indifferent: <score>} \\
\\ 
\\ 
\texttt{{[}End of answer{]}} \\
\\ 
\texttt{Remember: zero is a valid score, meaning they are likely not feeling that emotion. You must score at least one emotion > 0.} \\
\\ 
\texttt{Your answer:} \\
            }
        }
    }
     \\ \\

    \theutterance \stepcounter{utterance}  
    & \multicolumn{4}{p{0.6\linewidth}}{
        \cellcolor[rgb]{0.561,0.737,0.561}{
            \makecell[{{p{\linewidth}}}]{
                \texttt{\tiny{[P1$\rangle$GM]}}
                \texttt{Defeated: 0} \\
\texttt{Provoked: 10} \\
\texttt{Apologetic: 0} \\
\texttt{Indifferent: 0} \\
            }
        }
    }
    & & \\ \\

    \theutterance \stepcounter{utterance}  
    & & & \multicolumn{2}{p{0.3\linewidth}}{
        \cellcolor[rgb]{0.9,0.9,0.9}{
            \makecell[{{p{\linewidth}}}]{
                \texttt{\tiny{[GM$|$GM]}}
                \texttt{Defeated: 0} \\
\texttt{Provoked: 10} \\
\texttt{Apologetic: 0} \\
\texttt{Indifferent: 0} \\
            }
        }
    }
    & & \\ \\

    \theutterance \stepcounter{utterance}  
    & & & \multicolumn{2}{p{0.3\linewidth}}{
        \cellcolor[rgb]{0.9,0.9,0.9}{
            \makecell[{{p{\linewidth}}}]{
                \texttt{\tiny{[GM$|$GM]}}
                \texttt{Defeated: 0} \\
\texttt{Provoked: 7} \\
\texttt{Apologetic: 0} \\
\texttt{Indifferent: 0} \\
            }
        }
    }
    & & \\ \\

    \theutterance \stepcounter{utterance}  
    & & & \multicolumn{2}{p{0.3\linewidth}}{
        \cellcolor[rgb]{0.9,0.9,0.9}{
            \makecell[{{p{\linewidth}}}]{
                \texttt{\tiny{[GM$|$GM]}}
                \texttt{game\_result = LOSE} \\
            }
        }
    }
    & & \\ \\

    \theutterance \stepcounter{utterance}  
    & & \multicolumn{4}{p{0.6\linewidth}}{
        \cellcolor[rgb]{0.95,0.95,0.95}{
            \makecell[{{p{\linewidth}}}]{
                \texttt{\tiny{[P1$\leftrightarrow$P]}}
                \texttt{Post{-}episode Reflexion (not sent to GM):} \\
\texttt{New plan: The previous attempt failed because the scoring was too binary (10/0/0/0), likely missing the nuance that while Marcus is definitely provoked, the specific context of being called out on his own incompetence ("forgetting to show up") combined with his defensive "Low blow" response suggests a complex mix of emotions. Specifically, he feels **Provoked** (10) by the accusation, but he also feels **Defeated** (5) in the sense that his own lack of effort is being exposed and he is being publicly humiliated, forcing him to resort to aggression to regain control. "Apologetic" and "Indifferent" remain 0. I will adjust the scores to reflect this dual state of being attacked yet feeling personally defeated by the truth of his behavior.} \\
\\ 
\texttt{Defeated: 5} \\
\texttt{Provoked: 10} \\
\texttt{Apologetic: 0} \\
\texttt{Indifferent: 0} \\
            }
        }
    }
    & \\ \\

\end{supertabular}
}

\end{document}
